\documentclass{article}
\usepackage{tikz}
\usepackage{amsmath}
\usepackage[utf8]{inputenc}
\usepackage{xcolor}
\usepackage{ctex}  % 支持中文
\usepackage{geometry}

\geometry{a4paper, margin=1in}

\usetikzlibrary{shapes.geometric, arrows, positioning, fit, backgrounds, calc, shadows, chains, decorations.pathreplacing, matrix, shapes.symbols}

\begin{document}


\vspace{1cm}

\begin{tikzpicture}[
    node distance=1.5cm,
    box/.style={
        rectangle,
        draw=black,
        thick,
        fill=blue!10,
        text width=5cm,
        text centered,
        rounded corners,
        minimum height=1cm,
        drop shadow
    },
    smallbox/.style={
        rectangle,
        draw=black,
        thick,
        fill=blue!10,
        text width=3.8cm,
        text centered,
        rounded corners,
        minimum height=1.2cm,
        drop shadow
    },
    line/.style={
        draw,
        thick,
        -latex'
    },
    title/.style={
        font=\large\bfseries
    }
]

% 标题
\node[title] at (0, 0) (title2) {级联匹配策略};

% 级联匹配流程
\node[box, fill=green!10] at (0, -2) (classify) {根据轨迹质量进行分类:\\1. $\mathcal{T}_{high}$ - 高质量\\2. $\mathcal{T}_{medium}$ - 中等质量\\3. $\mathcal{T}_{low}$ - 低质量和新初始化轨迹};

% 级联处理
\node[box, below=of classify] (cascade) {逐级执行贪心匹配:\\按高质量→中等质量→低质量顺序处理\\确保稳定轨迹获得优先匹配权};

% 高质量轨迹匹配
\node[smallbox, fill=green!20, below=of cascade, xshift=-5cm] (high) {一\\匹配高质量轨迹\\确保稳定跟踪的目标\\先获得关联机会};

% 中等质量轨迹匹配
\node[smallbox, fill=yellow!20, below=of cascade] (medium) {二\\匹配中等质量轨迹\\处理暂时稳定的目标};

% 低质量轨迹匹配
\node[smallbox, fill=orange!20, below=of cascade, xshift=5cm] (low) {三\\匹配低质量轨迹\\处理新初始化或\\不稳定的目标跟踪};

% 动态阈值调整
\node[box, fill=blue!20, below=5.5cm of cascade] (dynamic) {动态阈值调整机制:\\$\tau_{IoU}(v,a) = \tau_0 \cdot e^{-\beta \cdot \|v\| \cdot a}$\\$\tau_0$: 基础阈值（0.7）\\$v$: 估计速度\\$a$: 目标面积\\$\beta$: 调整系数};


% 连接线
\draw [line] (classify) -- (cascade);
\draw [line] (cascade) -- (high);
\draw [line] (cascade) -- (medium);
\draw [line] (cascade) -- (low);

\draw [line] (high) |- ($(high.south)+(0,-1)$) -| (dynamic);
\draw [line] (medium) |- ($(medium.south)+(0,-0.5)$) -| (dynamic);
\draw [line] (low) |- ($(low.south)+(0,-1)$) -| (dynamic);

% 轨迹-检测匹配示意图
\begin{scope}[xshift=0cm, yshift=-18cm]
    \node[title] (title5) {轨迹-检测匹配示意图};

    % 轨迹
    \node[circle, draw, fill=green!20, minimum size=0.5cm] at (-2, -1.5) (T1) {$T_1$高};
    \node[circle, draw, fill=yellow!20, minimum size=0.5cm] at (-2, -3) (T2) {$T_2$中};
    \node[circle, draw, fill=orange!20, minimum size=0.5cm] at (-2, -4.5) (T3) {$T_3$低};

    % 检测
    \node[circle, draw, fill=blue!20, minimum size=0.5cm] at (2, -1.5) (D1) {$D_1$};
    \node[circle, draw, fill=blue!20, minimum size=0.5cm] at (2, -3) (D2) {$D_2$};
    \node[circle, draw, fill=blue!20, minimum size=0.5cm] at (2, -4.5) (D3) {$D_3$};

    % 连线
    \draw[thick, ->, green!50!black] (T1) -- node[above, sloped, font=\small] {先匹配} (D1);
    \draw[thick, ->, orange!50!black] (T2) -- node[above, sloped, font=\small] {次匹配} (D2);

    % 动态阈值
    \draw[thick, ->, dashed, red!50!black] (T3) -- node[above, sloped, font=\small] {$\tau_{动态}$} (D3);
\end{scope}

\end{tikzpicture}

\end{document}
